\section{Registered experiments and results}
\label{sec:results}

The manuscript never substitutes illustrative numbers for missing runs.  All
values in this section are TeX macros emitted by the checksum-validating
aggregator.  The present aggregate contains \NumCompletedRuns{} completed runs
and has status \emph{\ResultStatus}.

\subsection{Fidelity gate: reproducing the exact bridge}

Before interpreting any extension, we compare $M0/D0$ against the published
finite-$d$ reference curves under the same $L=2$, $\sigma=0.5$,
$\lambda=0.01$, summed loss, SGD step size, and informed basins.  The registered
agreement statistic is \ExactBridgeAgreement.  The gate requires: exact target
identity in unit tests; matching normalization and initial order parameters;
and curve agreement within combined finite-seed uncertainty.  Failure of this
gate blocks all architecture-level claims rather than being absorbed into a
new hyperparameter choice.

\subsection{Architecture robustness ($H1$--$H3$)}

The primary boundary estimate is \PrimaryBoundary{} with registered interval
\PrimaryBoundaryCI.  We compare boundary location and transition evidence
across tied/untied low-rank and QKVO models before adding heads.  For $M3$, the
phase label is supplemented by permutation-invariant head spectra so that a
head-index swap cannot masquerade as a discontinuity.  $M4$ pre/post-norm
experiments use matched seeds and residual interventions.  We report a shift
only if it survives learning-rate and optimizer controls.

\begin{figure}[t]
  \centering
  \IfFileExists{../assets/paper/architecture_ladder.png}{%
    \includegraphics[width=\linewidth]{../assets/paper/architecture_ladder.png}%
  }{\placeholderfigure{4.0cm}{Architecture-ladder results generated from the
  registered aggregate.}}
  \caption{\textbf{Architecture robustness of functional order.}  Panels share
  the same probe templates and hierarchical uncertainty.  Head-level panels
  are permutation invariant; missing state points remain marked rather than
  interpolated.}
  \label{fig:architecture-results}
\end{figure}

\subsection{Spectral chronology ($H4$)}

The registered difference between QK-outlier onset and semantic functional
onset is \OutlierLead{} optimizer steps, with interval \OutlierLeadCI.  This is
reported as precedence only if the interval excludes zero under independent
teacher resampling and remains stable to the predeclared outlier threshold.
Otherwise the correct conclusion is simultaneous or unresolved onset.

\subsection{Data ensembles and scaling paths ($H5$)}

The $D0$--$D5$ comparison tests whether a boundary is preserved after matching
second moments and then after deliberately breaking Gaussianity, continuity,
and token independence.  D4/D5 hidden states permit conditional analyses that
distinguish genuine semantic recovery from exploiting state persistence.

The preferred finite-size exponent is \FSSExponent{} with interval
\FSSExponentCI.  We do not combine $n\propto d$ and $n\propto d^2$ runs in one
collapse.  Competing paths are compared by held-out collapse score and
bootstrap stability over the same finite-width window.

\begin{figure}[t]
  \centering
  \IfFileExists{../assets/paper/fss_collapse.png}{%
    \includegraphics[width=\linewidth]{../assets/paper/fss_collapse.png}%
  }{\placeholderfigure{4.0cm}{Finite-size collapse, Binder crossing, and
  susceptibility panels generated from completed confirmatory runs.}}
  \caption{\textbf{Finite-size evidence.}  A non-analytic transition claim
  requires compatible Binder, susceptibility, distributional, and collapse
  evidence.  Smooth finite-width changes are described as crossovers.}
  \label{fig:fss-results}
\end{figure}

